\chapter{Conclusions}
\label{cap:cap5}

This project set out to address a recurring limitation in existing digital nutrition
platforms: the gap between tools that track what a user has already eaten and tools
that can proactively generate a nutritionally precise meal plan tailored to individual
metabolic targets. The result is a full-stack web application that automates daily
meal planning through a combination of machine learning-based meal filtering and
mathematical optimization, deployed as a cloud-native microservices system on
Microsoft Azure.

The core objective — generating a complete daily meal plan that satisfies a user's
calorie and macro-nutrient targets at the gram level — was achieved. The experimental
results presented in Chapter 4 confirm that the optimizer produces meal plans within
a 1.3\% deviation from the calculated calorie target across all tested activity groups
and dietary goals, which in absolute terms amounts to at most 32 kcal. This level of
precision exceeds what manual logging tools offer by design, as those platforms place
the burden of nutritional alignment entirely on the user.

The most significant personal contribution of this work is the two-stage optimization
pipeline at the core of the \texttt{OptimizerService}. Rather than relying on static
recipe lookups, the system first employs a Multi-Layer Perceptron classifier to
shortlist candidate meals based on the user's computed calorie and protein targets,
then applies a nonlinear constrained optimizer (SLSQP) to scale ingredient amounts
to gram-level precision by minimizing a weighted quadratic deviation from those targets. This combination — a learned classifier
followed by exact mathematical optimization — is not present in any of the comparable
platforms surveyed in Chapter 2.

\section{Comparison with Similar Solutions}

As outlined in the comparative analysis from Chapter 2 (see
Table~\ref{tab:comparative_analysis}), existing digital nutrition solutions fall into
two broad categories. Manual logging utilities such as MyFitnessPal require the user
to search and record food entries after consumption, offering no generative capability.
Dynamic metabolic trackers such as MacroFactor introduce adaptive calorie targets based
on observed weight trends, but still do not generate meal plans — the user remains
responsible for food selection. Standard automated planners can suggest recipes, but
treat them as fixed structures and cannot adjust ingredient quantities to meet a
specific calorie target.

The proposed system addresses all three limitations simultaneously: it generates a
complete plan automatically, selects meals intelligently through a trained classifier,
and scales each meal's ingredient amounts to satisfy the user's targets as hard
mathematical constraints. In this sense the application occupies a distinct position
in the landscape — closer to a personalized nutrition engine than a tracking tool.

\section{Future Development Directions}

The current implementation establishes a functional foundation, but several directions
could meaningfully extend the system's practical utility.

The most impactful near-term improvement would be a streamlined ingredient input
mechanism. At present, the optimizer works from a fixed read-only dataset of
ingredients bundled inside the container. A natural extension would be to allow users
to register the ingredients they currently have available — effectively describing the
contents of their fridge or pantry — and have the optimizer constrain its meal
generation to only those ingredients. This would shift the system from a general
meal planner to a context-aware one, reducing food waste and increasing the
likelihood that users actually prepare the suggested meals.

A second direction is real-world longitudinal validation. The experimental results in
Chapter 4 confirm numerical accuracy, but they do not measure whether consistent use
of the application over time produces the intended physiological outcomes. A structured
test — tracking a user's body weight, energy levels, and adherence rate over the course
of a month while following the generated plans — would provide evidence of the system's
practical effectiveness beyond algorithmic correctness. This kind of validation would
also surface usability issues that only emerge through daily use.

\section{Lessons Learned}

The microservices architecture adopted in this project introduced meaningful development
overhead — separate deployments, inter-service networking, CORS configuration, and
container registry management — that a monolithic design would not have required. In
retrospect, the split between the .NET API layer and the Python optimizer was
architecturally justified: the two services use different runtimes and have different
scaling characteristics, and a crash or redeployment of the optimizer does not
interrupt authentication or plan management. For projects of similar scope, a 
microservices boundary is most valuable when services differ in runtime, language, or 
failure tolerance — not simply to add architectural complexity.

Transitioning from a local Docker environment to Azure Container Apps required an
initial setup investment — configuring environment variables, internal service URLs,
and CORS origins correctly for the cloud environment. This friction is expected and
largely a one-time cost, but it reinforced that cloud deployment should be accounted
for early in the development cycle rather than treated as a final step.

The decision to use a nonlinear constrained optimizer for ingredient scaling proved
correct. An approximate approach — rounding ingredient amounts to the nearest fixed
increment — would have been simpler to implement but would have introduced uncontrolled
calorie deviation. The SLSQP solver finds the closest achievable solution within the
defined bounds and constraints, which is what makes the sub-1.3\% deviation result
meaningful.

Finally, bundling the SQLite dataset as a static artifact inside the container image
is a clean and dependency-free pattern for read-only data, but it does mean that any
change to the meal or ingredient dataset requires rebuilding and redeploying the
container. In practice this was not an issue, as the dataset was finalized before
deployment, but it is a real tradeoff to consider if the dataset were expected to
evolve over time.